\documentclass[addpoints]{exam}
\usepackage{amsmath}
\usepackage{amssymb}
\usepackage{color}
\usepackage{siunitx}

% \printanswers

\newcommand{\event}{Personal Finance}
\newcommand{\tournament}{}
\def\type{0}

\setlength\fillinlinelength{\textwidth}
\CorrectChoiceEmphasis{\color{red}\bfseries}
\SolutionEmphasis{\color{red}\bfseries}

\setlength\answerclearance{2pt}
\pagestyle{headandfoot}
\runningheadrule
\runningheader{\event}{\tournament}{\ifnum\type=1 Class Set Test \else Full Name:\hspace{1cm} \fi}
\runningfootrule
\runningfooter{}{Page \thepage\ of \numpages}{}

\begin{document}
\begin{center}
    {\huge Grade 11 Foundations for College Math \\ \event
    \ifprintanswers \space Key
    \else \space \ifnum\type<2 Test \fi \ifnum\type=1 \textbf{(CLASS SET)} \fi
          \ifnum\type=2 Answer Sheet \fi
    \fi \\ \vspace{3mm}\tournament\vspace{1mm}}
\end{center}

\vspace{.25\textheight}

\begin{center}
    First Name: \ifprintanswers
    \underline{\hspace{3.5cm}}\textbf{KEY}\underline{\hspace{5.5cm}}\\\ \vspace{5mm}
    \else \underline{\hspace{10cm}}\\ \vspace{5mm} \fi
    Last Name: \underline{\hspace{10cm}}\\ \vspace{5mm}
\end{center}

\noindent\textbf{\underline{Directions:}}
\begin{itemize}
    \ifnum\type=1 \item \textbf{This test is a class set.} Do not write on this paper. \fi
    \item Show all work for full marks. Round dollar amounts to the nearest cent.
    \item The test is designed to be completed in 75 minutes.
\end{itemize}
\begin{center}
\textit{For grading use only}\\
\multirowgradetable{1}[pages]
\end{center}

\newpage
\begin{questions}

\section*{Part A --- Multiple Choice (15 marks)}
\vspace{5mm}

% Q1
\question[1]
Simple interest is calculated using the formula $I = Prt$. If $P = \$1{,}200$, $r = 0.05$, and $t = 3$ years, what is the interest earned?
\begin{choices}
    \choice \$60.00
    \correctchoice \$180.00
    \choice \$216.00
    \choice \$1{,}380.00
\end{choices}

% Q2
\question[1]
Which formula gives the accumulated amount $A$ for compound interest?
\begin{choices}
    \choice $A = P + Prt$
    \correctchoice $A = P(1+i)^n$
    \choice $A = P \cdot r \cdot t$
    \choice $A = P(1-i)^n$
\end{choices}

% Q3
\question[1]
\$500 is invested at 6\% per year compounded annually for 2 years. What is the accumulated amount?
\begin{choices}
    \choice \$560.00
    \correctchoice \$561.80
    \choice \$600.00
    \choice \$530.00
\end{choices}

% Q4
\question[1]
The nominal annual interest rate is 8\% compounded quarterly. What is the interest rate per compounding period?
\begin{choices}
    \choice 8\%
    \choice 4\%
    \correctchoice 2\%
    \choice 1\%
\end{choices}

% Q5
\question[1]
An investment compounds semi-annually. If the annual rate is 6\%, how many compounding periods are there in 3 years?
\begin{choices}
    \choice 3
    \choice 4
    \correctchoice 6
    \choice 12
\end{choices}

% Q6
\question[1]
Which statement correctly compares simple and compound interest over time?
\begin{choices}
    \choice Simple interest always earns more than compound interest.
    \correctchoice Compound interest earns more than simple interest over the same period at the same rate.
    \choice They always produce the same final amount.
    \choice Simple interest is better for long-term investments.
\end{choices}

% Q7
\question[1]
Present value is best described as:
\begin{choices}
    \choice The interest earned on an investment.
    \choice The future accumulated amount.
    \correctchoice The amount that must be invested today to reach a desired future value.
    \choice The nominal interest rate.
\end{choices}

% Q8
\question[1]
How much must you invest today at 5\% compounded annually to have \$1{,}050 in one year?
\begin{choices}
    \correctchoice \$1{,}000.00
    \choice \$997.62
    \choice \$1{,}050.00
    \choice \$952.38
\end{choices}

% Q9
\question[1]
A credit card charges 19.99\% per year compounded daily. What is the approximate daily interest rate?
\begin{choices}
    \choice 0.1999\%
    \correctchoice 0.0548\%
    \choice 1.666\%
    \choice 0.5479\%
\end{choices}

% Q10
\question[1]
Which of the following describes a GIC?
\begin{choices}
    \choice A loan product with variable interest.
    \correctchoice A deposit investment that earns a guaranteed fixed interest rate for a set term.
    \choice A type of mutual fund.
    \choice A government-issued bond with no fixed return.
\end{choices}

% Q11
\question[1]
A mortgage is best described as:
\begin{choices}
    \choice A short-term personal loan.
    \correctchoice A long-term loan secured by real estate, repaid in regular installments.
    \choice A credit card with a high limit.
    \choice An investment account.
\end{choices}

% Q12
\question[1]
\$2{,}000 is owed on a credit card at 20\% per year compounded monthly. If no payment is made, approximately how much interest accrues in one month?
\begin{choices}
    \choice \$40.00
    \correctchoice \$33.33
    \choice \$400.00
    \choice \$3.33
\end{choices}

% Q13
\question[1]
An investment of \$1{,}000 at 4\% compounded annually for 10 years gives approximately:
\begin{choices}
    \choice \$1{,}400.00
    \correctchoice \$1{,}480.24
    \choice \$1{,}040.00
    \choice \$1{,}500.00
\end{choices}

% Q14
\question[1]
If you make only the minimum payment on a credit card balance each month, what happens over time?
\begin{choices}
    \choice The balance decreases quickly.
    \choice The balance stays the same.
    \correctchoice The balance decreases very slowly and you pay a lot of interest.
    \choice The interest rate decreases.
\end{choices}

% Q15
\question[1]
Which compounding frequency results in the highest accumulated value for the same principal, rate, and time?
\begin{choices}
    \choice Annually
    \choice Semi-annually
    \choice Monthly
    \correctchoice Daily
\end{choices}

\section*{Part B --- Short Answer (33 marks)}

% Q16
\question[6]
\textbf{Simple Interest.} Use the formula $I = Prt$ and $A = P + I$.

\begin{parts}
    \part[2] Kenji invests \$3{,}500 at a simple interest rate of 4.5\% per year for 2 years. How much interest does he earn, and what is the accumulated amount?
    \begin{solution}
        $I = Prt = 3500 \times 0.045 \times 2 = \$315.00$\\
        $A = P + I = 3500 + 315 = \$3{,}815.00$
    \end{solution}

    \part[2] Maria needs \$1{,}218.75 and invests \$1{,}125 at a simple interest rate of $r$ per year for 3 years. Find the annual interest rate.
    \begin{solution}
        $I = A - P = 1218.75 - 1125 = \$93.75$\\
        $I = Prt \Rightarrow 93.75 = 1125 \times r \times 3$\\
        $r = \dfrac{93.75}{3375} = 0.0278 = 2.78\%$
    \end{solution}

    \part[2] An investment of \$2{,}000 at 6\% simple interest earns \$480. For how many years was the money invested?
    \begin{solution}
        $I = Prt \Rightarrow 480 = 2000 \times 0.06 \times t$\\
        $480 = 120t$\\
        $t = 4 \text{ years}$
    \end{solution}
\end{parts}

% Q17
\question[6]
\textbf{Compound Interest.}

\begin{parts}
    \part[2] Find the accumulated amount if \$4{,}000 is invested at 6\% per year compounded quarterly for 5 years.
    \begin{solution}
        $i = \dfrac{0.06}{4} = 0.015$, \quad $n = 5 \times 4 = 20$\\
        $A = P(1+i)^n = 4000(1.015)^{20}$\\
        $(1.015)^{20} \approx 1.34686$\\
        $A \approx 4000 \times 1.34686 = \$5{,}387.42$
    \end{solution}

    \part[2] The same \$4{,}000 is invested at 6\% per year compounded semi-annually for 5 years. Find the accumulated amount.
    \begin{solution}
        $i = \dfrac{0.06}{2} = 0.03$, \quad $n = 5 \times 2 = 10$\\
        $A = 4000(1.03)^{10}$\\
        $(1.03)^{10} \approx 1.34392$\\
        $A \approx 4000 \times 1.34392 = \$5{,}375.67$
    \end{solution}

    \part[2] How does each compound interest result in parts (a) and (b) compare to simple interest at 6\% per year for 5 years on the same \$4{,}000? Which compounding option is best?
    \begin{solution}
        Simple interest: $I = 4000 \times 0.06 \times 5 = \$1{,}200$, so $A = \$5{,}200.00$\\
        Quarterly compounding gives \$5{,}387.42 --- highest.\\
        Semi-annual compounding gives \$5{,}375.67.\\
        Both compound options outperform simple interest. Quarterly compounding is best because the interest is added to the principal more frequently.
    \end{solution}
\end{parts}

% Q18
\question[5]
\textbf{Present Value.} How much money must be invested today at 4\% per year compounded annually in order to have \$10{,}000 in 5 years? Round to the nearest cent.
\begin{solution}
    Use $A = P(1+i)^n$ and solve for $P$:
    \[P = \frac{A}{(1+i)^n} = \frac{10000}{(1.04)^5}\]
    $(1.04)^5 \approx 1.21665$
    \[P = \frac{10000}{1.21665} \approx \$8{,}219.27\]
    You must invest \$8{,}219.27 today.
\end{solution}

% Q19
\question[8]
\textbf{Credit Card Interest.} Samantha has a credit card balance of \$2{,}000. The card charges 19.99\% per year compounded daily. Assume 1 month = 30 days.

\begin{parts}
    \part[3] Find the balance after 30 days if no payment is made.
    \begin{solution}
        Daily rate: $i = \dfrac{0.1999}{365} \approx 0.00054767$\\
        Number of periods: $n = 30$\\
        $A = 2000(1 + 0.00054767)^{30}$\\
        $(1.00054767)^{30} \approx 1.01649$\\
        $A \approx 2000 \times 1.01649 = \$2{,}032.98$\\
        Interest charged: $\$2{,}032.98 - \$2{,}000.00 = \$32.98$
    \end{solution}

    \part[3] If instead Samantha pays off the full \$2{,}000 balance immediately (before any interest accrues), how much does she save compared to making no payment?
    \begin{solution}
        If she pays in full immediately, she owes exactly \$2{,}000 --- no interest.\\
        By not paying, she owes \$2{,}032.98 after 30 days.\\
        Savings $= \$2{,}032.98 - \$2{,}000.00 = \$32.98$\\
        She saves \$32.98 by paying in full right away.
    \end{solution}

    \part[2] Explain in 2--3 sentences why making only the minimum monthly payment on a credit card is financially costly over the long term.
    \begin{solution}
        When only the minimum payment is made, the remaining balance continues to accumulate interest at a high rate (nearly 20\% per year). Over time the total interest paid can far exceed the original purchase price. This means it can take many years and hundreds of extra dollars to pay off a relatively small balance.
    \end{solution}
\end{parts}

% Q20
\question[8]
\textbf{GIC Investment Comparison.} David has \$5{,}000 to invest for 3 years. He is comparing two options:
\begin{itemize}
    \item \textbf{Option A:} A 3-year GIC at 3.5\% per year compounded semi-annually.
    \item \textbf{Option B:} A savings account at 3.2\% per year compounded monthly.
\end{itemize}

\begin{parts}
    \part[3] Find the accumulated amount for Option A after 3 years.
    \begin{solution}
        $i = \dfrac{0.035}{2} = 0.0175$, \quad $n = 3 \times 2 = 6$\\
        $A = 5000(1.0175)^{6}$\\
        $(1.0175)^{6} \approx 1.10969$\\
        $A \approx 5000 \times 1.10969 = \$5{,}548.44$
    \end{solution}

    \part[3] Find the accumulated amount for Option B after 3 years.
    \begin{solution}
        $i = \dfrac{0.032}{12} \approx 0.0026\overline{6}$, \quad $n = 3 \times 12 = 36$\\
        $A = 5000(1.002\overline{6})^{36}$\\
        $(1.002\overline{6})^{36} \approx 1.10023$\\
        $A \approx 5000 \times 1.10023 = \$5{,}501.17$
    \end{solution}

    \part[2] Which option should David choose, and why? Mention at least one advantage and one disadvantage of the GIC compared to the savings account.
    \begin{solution}
        Option A (GIC) gives \$5{,}548.44 vs.\ Option B (savings) \$5{,}501.17, so the GIC earns \$47.27 more.\\
        \textbf{Advantage of GIC:} Guaranteed, locked-in rate --- David knows exactly what he will earn.\\
        \textbf{Disadvantage of GIC:} The money is locked in for 3 years; early withdrawal may result in a penalty or loss of interest.\\
        David should choose the GIC if he does not need access to the funds during the 3-year term.
    \end{solution}
\end{parts}

\end{questions}
\end{document}
